\section{Problem and Experimental Setting}
\label{sec:setting}

\paragraph{Synthetic worlds.} Each world seed generates \res{ncls} Gaussian classes in $\mathbb{R}^{\res{dim}}$
and several shifted domains, each an independent per-feature affine map of the source distribution plus small
noise. Disjoint splits are drawn from the same seed: labelled source data (to train the classifier), meta
data under $4$ simulated shifts (used only to fit subspaces), development and calibration streams with a
development holdout under $2$ validation shifts, and a test stream and test holdout under
\res{ntestdom} test shifts. The test stream has \res{nstream} samples and the test holdout
\res{nhold}. No test label is visible to any adaptation arm (Appendix~\ref{app:access}).

\paragraph{Model and entropy-minimization toy baseline.} A linear softmax classifier $f(x)=W z+b$ is trained
on source data and frozen. Test-time adaptation updates $\theta=(\gamma,\beta)\in\mathbb{R}^{2\times\res{dim}}$
in $z=\gamma\odot\hat x+\beta$, where $\hat x$ is standardized with \emph{fixed source statistics}; there are
no running buffers and no optimizer state. Each step takes one plain-SGD step, with constant learning rate
$\eta$, on the mean prediction entropy of the current batch. We call this the \emph{entropy-minimization toy
baseline} (EM-toy). It shares Tent's objective but differs from the published Tent configuration, which
adapts normalization affine parameters of a deep network and re-estimates normalization statistics on each
test batch~\citep{wang2021tent}. Run identifiers keep the historical arm id \texttt{tent\_full}; the
mapping to display names is in \path{paper/tables/display_names.tsv}.

\paragraph{Orders and conditions.} One multiset of test sample IDs is replayed in \res{norders} fixed orders
per world seed, and every arm sees the same orders. In the \emph{primary stationary} condition the stream is
partitioned once into batches of \res{batch} and orders permute only the batches (\res{stepsstat} steps),
so batch composition is identical across orders. The \emph{secondary} condition reshuffles samples and
re-forms batches for every order; it is reported separately and not used for the decision. In the
\emph{drift} condition each test domain is partitioned once into single-domain batches and orders permute the
sequence of domain blocks and the batches within each block (\res{stepsdrift} steps). Before every run the
harness checks that the adapter starts from the source parameters and $\theta_0=(\mathbf{1},\mathbf{0})$,
that frozen parameters are unchanged afterwards, and that the global random state was not consumed.

\paragraph{Metrics.} For the stationary conditions the primary metric is the terminal error on a holdout that
is shared by all orders and arms. For drift we do not score the mixture holdout, which would penalise
adaptation to the latest regime. Instead we score the current regime's holdout at the end of every domain
block and average over blocks, and we report predict-then-update online error separately. For each world seed
and arm we report the mean over the observed orders, the standard deviation (SD) across them, the worst
observed order and the gain over no adaptation. Orders are nested in world seeds. We therefore report
per-seed values and paired per-seed differences (candidate minus control over the same orders), and we form
no confidence interval that pools seeds and orders. With \res{nseeds} seeds, the independent units are the
seeds.

\paragraph{Arms and information access.} All adapting arms use the same loss, the same number of steps and no
test labels.
\begin{itemize}
\item \emph{No adaptation}: the frozen source model.
\item \emph{EM-toy (full affine)}: all $2\times\res{dim}$ parameters, with a tuned $\eta$.
\item \emph{EM-toy, lr$\times$0.3 and lr$\times$0.1}: small-step controls.
\item \emph{Order-penalized subspace (candidate)}: Section~\ref{sec:analysis}; it uses labelled meta data at fit
time.
\item \emph{Utility-only subspace}: the same meta batches, meta labels, candidate pool and second-order
statistics as the candidate, restricted to the candidate's realised rank, with the order penalty set to zero.
\item \emph{Per-step norm-matched (causal)}: full-space direction, with step $t$ rescaled to the norm of step
$t$ of a lockstep replica of the candidate, so no future step is readable.
\item \emph{Global norm-matched}: a full-space $\eta$ bisected on the calibration stream so that the mean
terminal displacement matches the candidate's; this reads no labels.
\end{itemize}
The EM-toy, candidate and utility-only arms select $\eta$ from the same grid \res{lrgrid} on development
streams, by development holdout error (mean current-regime error for drift). Ties go to the smaller value.
The grid was not widened for this study.

\paragraph{Pre-registered decision rule.} The rule, the seeds (the first three of five pre-registered world
seeds and the first eight of sixteen orders) and the minimum effects of interest (MIE; gain
\res{mieg}\,pp, SD \res{mies}\,pp, tolerance \res{tolg}\,pp) were committed before the run (commit
\texttt{\res{commit}}). The candidate is flagged for a later real-data pilot only if all of the following
hold:
\begin{enumerate}
\item In the primary stationary condition, its gain over no adaptation is at least the MIE on average and
positive in every seed.
\item Its SD reduction relative to EM-toy is at least the MIE on average and positive in every seed.
\item It \emph{dominates} each required control. That is, it is better by the MIE in mean error, or in SD,
with the same sign in every seed, and not worse than the tolerance on the other axis.
\item In drift, its current-regime and online errors are within tolerance of the utility-only control and
not worse than no adaptation.
\end{enumerate}
A required cell that failed, hit a cap or did not run makes the result \emph{incomplete}; otherwise the result
is \emph{not supported}.
